\documentclass[11pt,a4paper]{article}
\usepackage[utf8]{inputenc}
\usepackage[T1]{fontenc}
\usepackage[spanish]{babel}
\usepackage{hyperref}
\usepackage{geometry}
\geometry{margin=2.5cm}
\usepackage{parskip}

\title{111 el candado no nombra series}
\author{Alejandro Valenzuela Rivera\\\small Universidad Autónoma de Nuevo León}
\date{\today}
\begin{document}
\maketitle

\begin{quote}
\textbf{Sesgo declarado.} Este trabajo se sostiene desde el sesgo declarado de su autor, y no finge neutralidad. Lo que afirma se puede tumbar, y ese es el punto: el falsificador está escrito, no escondido.
\end{quote}

\section*{Resumen}
Una tesis que se sostiene sola no necesita apoyarse en un personaje para que se entienda. La escena que enseno algo se cuenta como escena: lo que paso, lo que se entendio, y cuando. Sin el nombre de la obra, porque el nombre no sostiene el argumento.

\section{Por qué se sostiene}
Se lee raro porque el nombre de la obra arrastra el tono de la obra: parece que se argumenta con una caricatura en vez de con una idea. Y el argumento pierde cuando el lector se distrae pensando en la serie. Neutral no es frio: es que la idea se sostenga sin el muleta.

\section{De dónde viene}
La autoridad, 2026-10-26.

\section{El caso}
Ninguna: es la regla del candado.

\section{Método}
Un candado que nombra lo que pregunta, sin citar la serie como autoridad.

\section{Límites, y qué lo tumba}
Un candado que se sostiene porque la serie dice X.

\section{Lo que queda abierto}
\textquestiondown{}Si quito el nombre de la obra, la idea sigue de pie?

\section*{Cómo se leyó esta tesis}
Esto se escribió en el repositorio man-made-empathy, del que este candado es una pieza: candados/03 el nombre y la obra/111 el candado no nombra series.md. La máquina que lo compuso lo hizo desde este mismo archivo, y no agregó nada que no estuviera escrito. Si algo de esto es falso, se dice en el mismo sitio, y se corrige con fecha.

\end{document}
